\chapter{Witness Presheaf Value Tables}
\label{appC:chapter}

This appendix tabulates the two witness presheaves $F_A$ and $F_R$ in
full: their value sets at all sixteen contexts, as counts and as explicit
bit-assignment sets; their restriction maps; and the gluing computations
at the global context $\Omega$ for both topologies, which are the
registered core of the distinctness theorem (Theorem NT2-3, Chapter~9)
and of the
certificate verdicts
\texttt{acbe\_witness\_is\_acbe\_sheaf} /
\texttt{acbe\_witness\_is\_rpm\_sheaf} and their RPM mirrors
(Chapter~11). Every entry is derived directly from the definitions; the
derivations are shown. Context names and complete-block containments are
those of Appendix~A.

\section{Definitions and the counting rule}
\label{appC:sec:defs}

\begin{notation}
\label{appC:not:witness}
$F_A(V)$ is the set of functions from
$\{\, i \in \{0,1\} \mid A_i \subseteq V \,\}$ to $\{0, 1\}$: bit
assignments on the \emph{complete} ACBE blocks contained in $V$. For
$U \subseteq V$, the restriction map forgets the coordinates $i$ with
$A_i \subseteq V$ but $A_i \not\subseteq U$ and keeps the remaining bits
unchanged. $F_R$ is defined symmetrically on complete RPM blocks. We
write a section of $F_A(V)$ as a tuple over its index set: at $\Omega$
the four sections are $(b_0, b_1)$ with $b_i \in \{0,1\}$ the bit on
$A_i$; at a context with one complete block the two sections are $(0)$
and $(1)$; at a context with no complete block the unique section is the
empty assignment, written $()$.
\end{notation}

The counting rule is immediate: $|F_A(V)| = 2^{\alpha(V)}$ where
$\alpha(V)$ is the number of complete $A$-blocks contained in $V$
(Definition of complete blocks, Chapter~7), and likewise
$|F_R(V)| = 2^{\rho(V)}$. The containment table of Appendix~A gives
$\alpha(V)$ and $\rho(V)$ at every context. Two facts drive every verdict
below:

\begin{finding}
\label{appC:fnd:novaluecrossing}
No RPM piece ever contains a complete ACBE block: for every context $V$
and every $j$, $V \cap R_j \subseteq R_j$, and $A_i \subseteq R_j$ is
impossible since $A_i \cap R_j = \{(i,j)\}$ is a proper subset of
$A_i$. Hence $F_A(V \cap R_j) = \{()\}$ is a one-element set for every
$V$ and $j$. Symmetrically, $F_R(V \cap A_i)$ is a one-element set for
every $V$ and $i$. This single observation determines the entire
failure pattern of the cross gluing computations below.
\end{finding}

\section{Value tables}
\label{appC:sec:values}

\begin{longtable}{@{}lccll@{}}
\toprule
Context $V$ & $|F_A(V)|$ & $|F_R(V)|$ & $F_A(V)$ explicitly &
$F_R(V)$ explicitly \\
\midrule
\endhead
\bottomrule
\endfoot
$\varnothing$ & 1 & 1 & $\{()\}$ & $\{()\}$ \\
$\{\omega_{00}\}$ & 1 & 1 & $\{()\}$ & $\{()\}$ \\
$\{\omega_{01}\}$ & 1 & 1 & $\{()\}$ & $\{()\}$ \\
$\{\omega_{10}\}$ & 1 & 1 & $\{()\}$ & $\{()\}$ \\
$\{\omega_{11}\}$ & 1 & 1 & $\{()\}$ & $\{()\}$ \\
$A_0$ & 2 & 1 & $\{(0), (1)\}$ on $A_0$ & $\{()\}$ \\
$A_1$ & 2 & 1 & $\{(0), (1)\}$ on $A_1$ & $\{()\}$ \\
$R_0$ & 1 & 2 & $\{()\}$ & $\{(0), (1)\}$ on $R_0$ \\
$R_1$ & 1 & 2 & $\{()\}$ & $\{(0), (1)\}$ on $R_1$ \\
$\Delta$ & 1 & 1 & $\{()\}$ & $\{()\}$ \\
$\nabla$ & 1 & 1 & $\{()\}$ & $\{()\}$ \\
$T_{00}$ & 2 & 2 & $\{(0), (1)\}$ on $A_1$ & $\{(0), (1)\}$ on $R_1$ \\
$T_{01}$ & 2 & 2 & $\{(0), (1)\}$ on $A_1$ & $\{(0), (1)\}$ on $R_0$ \\
$T_{10}$ & 2 & 2 & $\{(0), (1)\}$ on $A_0$ & $\{(0), (1)\}$ on $R_1$ \\
$T_{11}$ & 2 & 2 & $\{(0), (1)\}$ on $A_0$ & $\{(0), (1)\}$ on $R_0$ \\
$\Omega$ & 4 & 4 & $\{(b_0, b_1) : b_i \in \{0,1\}\}$ &
$\{(c_0, c_1) : c_j \in \{0,1\}\}$ \\
\end{longtable}

Derivation: each row applies the counting rule to the complete-block
columns of Appendix~A. Nine contexts contain no complete $A$-block
(value $1$), six contain exactly one ($A_0$, $A_1$, and the four
triples; value $2$), and $\Omega$ contains both (value $4$); mirrored
for $F_R$. The section totals over all contexts are
$9 \cdot 1 + 6 \cdot 2 + 4 = 25$ for each witness.

\section{Restriction maps}
\label{appC:sec:restrictions}

Restriction in $F_A$ is determined by which complete blocks survive the
inclusion. The following table describes every qualitatively distinct
class of inclusions; all others are composites or have singleton source
or target value sets (where the map is the unique one). Bits are named
as in the value table.

\begin{longtable}{@{}llll@{}}
\toprule
Inclusion class & Surviving $A$-blocks & $F_A$ restriction map &
Character \\
\midrule
\endhead
\bottomrule
\endfoot
$\Omega \supseteq T_{0j}$ & $A_1$ &
$(b_0, b_1) \mapsto (b_1)$ & drops $b_0$; two-to-one \\
$\Omega \supseteq T_{1j}$ & $A_0$ &
$(b_0, b_1) \mapsto (b_0)$ & drops $b_1$; two-to-one \\
$\Omega \supseteq A_i$ & $A_i$ &
$(b_0, b_1) \mapsto (b_i)$ & drops the other bit; two-to-one \\
$\Omega \supseteq R_j, \Delta, \nabla$, singletons, $\varnothing$ &
none & $(b_0, b_1) \mapsto ()$ & constant; four-to-one \\
$T_{ij} \supseteq A_{1-i}$ & $A_{1-i}$ &
$(b) \mapsto (b)$ & identity on the surviving bit \\
$T_{ij} \supseteq$ any $U$ with $A_{1-i} \not\subseteq U$ & none &
$(b) \mapsto ()$ & constant; two-to-one \\
$A_i \supseteq$ proper $U$ & none &
$(b) \mapsto ()$ & constant; two-to-one \\
source with $k_A = 0$ & none & $() \mapsto ()$ & identity of a
singleton \\
\end{longtable}

Functoriality is visible in the table: a composite of inclusions
survives a bit exactly when both steps do, and on surviving bits every
map is the identity, so both functor laws hold coordinatewise. This is
the content the certificate validates as
\texttt{acbe\_witness\_functorial} before computing any verdict; the
$F_R$ table is the mirror image under the coordinate exchange (replace
$A$-blocks by $R$-blocks, $T_{ij} \supseteq A_{1-i}$ by
$T_{ij} \supseteq R_{1-j}$, and $b_i$ by $c_j$), validated as
\texttt{rpm\_witness\_functorial}.

\section{Gluing at $\Omega$: the four registered computations}
\label{appC:sec:gluing}

The sheaf condition at $\Omega$ compares the value set with the set of
matching families over each global cover. Since distinct pieces of
either cover meet in $\varnothing$, where both witnesses have the
one-element value $\{()\}$, compatibility is automatic and the
matching-family set is the plain product of the piece values. The four
computations, tabulated:

\begin{longtable}{@{}llccl@{}}
\toprule
Witness / cover & Piece values & Matching families & Value at $\Omega$ &
Gluing map and verdict \\
\midrule
\endhead
\bottomrule
\endfoot
$F_A$ over $K_A(\Omega)$ & $F_A(A_0) = \{(0),(1)\}$,
$F_A(A_1) = \{(0),(1)\}$ & $2 \times 2 = 4$ & 4 &
$(b_0, b_1) \mapsto ((b_0), (b_1))$: bijection --- sheaf condition
holds \\
$F_A$ over $K_R(\Omega)$ & $F_A(R_0) = \{()\}$, $F_A(R_1) = \{()\}$ &
$1 \times 1 = 1$ & 4 &
all four of $(0,0), (0,1), (1,0), (1,1) \mapsto ((),())$: not
injective --- fails \\
$F_R$ over $K_R(\Omega)$ & $F_R(R_0) = \{(0),(1)\}$,
$F_R(R_1) = \{(0),(1)\}$ & $2 \times 2 = 4$ & 4 &
$(c_0, c_1) \mapsto ((c_0), (c_1))$: bijection --- sheaf condition
holds \\
$F_R$ over $K_A(\Omega)$ & $F_R(A_0) = \{()\}$, $F_R(A_1) = \{()\}$ &
$1 \times 1 = 1$ & 4 &
four-to-one onto $((),())$: not injective --- fails \\
\end{longtable}

Derivations. Row 1: a matching family for $K_A(\Omega)$ is a free pair
(one section over each $A_i$; the overlap $A_0 \cap A_1 = \varnothing$
imposes agreement in a singleton), and the canonical map sends the
global assignment $(b_0, b_1)$ to its restrictions, which by the
restriction table are $(b_0)$ on $A_0$ and $(b_1)$ on $A_1$; this is
the evident bijection, so $F_A$ satisfies the ACBE gluing condition at
$\Omega$. Row 2: by Finding~\ref{appC:fnd:novaluecrossing} each RPM
half carries the one-element value, so exactly one matching family
exists; the canonical map sends each of the four global sections to it,
failing injectivity, so $F_A$ is not a $J_R$-sheaf --- this is the
registered ``$4 \to 1$'' computation, exhibited here element by
element. Rows 3 and 4 are the mirror computations.

\section{Gluing at the remaining contexts}
\label{appC:sec:othercontexts}

The certificate's verdicts quantify over \emph{all} contexts, not only
$\Omega$; for completeness we tabulate the matching-family counts
everywhere. For the own-partition cover the count is
$|F_A(V \cap A_0)| \cdot |F_A(V \cap A_1)|$ (empty pieces contribute
the factor $1$); for the cross cover it is identically $1$ by
Finding~\ref{appC:fnd:novaluecrossing}. At $\varnothing$ the null
cover has exactly one matching family.

\begin{longtable}{@{}lcccc@{}}
\toprule
Context $V$ & $|F_A(V)|$ & Matching for $K_A(V)$ & Matching for
$K_R(V)$ & Verdicts for $F_A$ (own / cross) \\
\midrule
\endhead
\bottomrule
\endfoot
$\varnothing$ & 1 & 1 & 1 & holds / holds \\
$\{\omega_{ij}\}$ (each) & 1 & $1 \cdot 1 = 1$ & 1 & holds / holds \\
$A_0$ & 2 & $2 \cdot 1 = 2$ & 1 & holds / fails ($2 \to 1$) \\
$A_1$ & 2 & $1 \cdot 2 = 2$ & 1 & holds / fails ($2 \to 1$) \\
$R_0$ & 1 & $1 \cdot 1 = 1$ & 1 & holds / holds \\
$R_1$ & 1 & $1 \cdot 1 = 1$ & 1 & holds / holds \\
$\Delta$ & 1 & $1 \cdot 1 = 1$ & 1 & holds / holds \\
$\nabla$ & 1 & $1 \cdot 1 = 1$ & 1 & holds / holds \\
$T_{00}$ & 2 & $1 \cdot 2 = 2$ & 1 & holds / fails ($2 \to 1$) \\
$T_{01}$ & 2 & $1 \cdot 2 = 2$ & 1 & holds / fails ($2 \to 1$) \\
$T_{10}$ & 2 & $2 \cdot 1 = 2$ & 1 & holds / fails ($2 \to 1$) \\
$T_{11}$ & 2 & $2 \cdot 1 = 2$ & 1 & holds / fails ($2 \to 1$) \\
$\Omega$ & 4 & $2 \cdot 2 = 4$ & 1 & holds / fails ($4 \to 1$) \\
\end{longtable}

Reading the table: in the own-partition column the matching count
equals the value count at every context, and inspection of the
canonical maps (as in Row~1 of the $\Omega$ table; at contexts with an
identity piece the condition is trivial) shows each is a bijection ---
$F_A$ is a $J_A$-sheaf, and by the basis criterion the generating
families suffice. In the cross column the matching count is constantly
$1$, so the gluing map fails injectivity at exactly the seven contexts
where $|F_A(V)| \ge 2$: the two $A$-blocks, the four triples, and
$\Omega$. The registered $\Omega$ computation is thus one instance of
a uniform failure pattern; any one of the seven already establishes
that $F_A$ is not a $J_R$-sheaf, and the certificate's exhaustive run
confirms the pattern at all of them. The mirrored table for $F_R$
(exchange the two partitions, the two block families, and the two
cover columns) shows $F_R$ is a $J_R$-sheaf failing $J_A$-gluing at
the two $R$-blocks, the four triples, and $\Omega$.

\section{Reflector values on the witnesses}
\label{appC:sec:reflectorvalues}

For completeness we tabulate what the three reflectors of Chapter~10 do
to the witnesses of Chapter~9, since these values are quoted in the body (the
annihilation example and the both-routes example). Each entry is
computed from the reflector formula
$(R_P F)(V) = \prod_{B \in I_P(V)} F(V \cap B)$ together with
Finding~\ref{appC:fnd:novaluecrossing} and the value tables above.

\begin{longtable}{@{}lccc@{}}
\toprule
Presheaf & $R_A$ applied & $R_R$ applied & $R_Q$ applied \\
\midrule
\endhead
\bottomrule
\endfoot
$F_A$ & $\cong F_A$ (unit iso; $F_A$ is a $J_A$-sheaf) &
terminal ($1$ at every context) & terminal \\
$F_R$ & terminal ($1$ at every context) &
$\cong F_R$ (unit iso) & terminal \\
\end{longtable}

Derivations. $R_A F_A \cong F_A$: the unit of a reflector is an
isomorphism exactly on sheaves (Chapter~10), and the own-partition
column of Section~\ref{appC:sec:othercontexts} verifies the sheaf
condition. $R_A F_R$ terminal: every piece $V \cap A_i$ lies in a
single ACBE block, hence contains no complete RPM block, so every
factor $F_R(V \cap A_i)$ is a one-element set and the product is a
one-element set at every context. $R_Q F_A$ terminal: every factor
$F_A(\{\omega\})$ is a one-element set since a singleton contains no
complete block. The terminal presheaf is a sheaf for all three
topologies (every matching-family set of a constant-singleton presheaf
is a singleton, including at the null cover), so these values are
consistent with the sheaf lemma. The table also exhibits, in minimal
form, the content of the composition theorem on these fixtures: both
composite routes applied to either witness give the terminal presheaf,
so the two routes agree here for the strongest possible reason ---
they agree with each other by agreeing with $R_Q$.

\section{Section counts at a glance}
\label{appC:sec:glance}

Summing the value tables: each witness has $25$ sections in total over
the sixteen contexts ($9$ contexts contributing one section, $6$
contributing two, $\Omega$ contributing four); each reflected
cross-image ($R_A F_R$, $R_R F_A$, $R_Q F_A$, $R_Q F_R$) has exactly
$16$ sections, one per context. The asymmetry $25$ versus $16$ is the
quantitative footprint of nontriviality: a modality that were
naturally isomorphic to the identity would preserve section counts at
every context, and the drop from $(1,2,4)$-valued to constantly
$1$-valued is what the certificate's false verdicts measure
context-by-context.

\begin{boundary}
\label{appC:bnd:scope}
These tables establish combinatorial facts about two specific finite
presheaves on $L$ and nothing more. They corroborate the distinctness
theorem's witnesses (Theorem NT2-3, Chapter~9, scoped to
$\PSh(L, \Set)$); they
attach no interpretation to the ACBE and RPM marks, and they license
no ambient, probabilistic, empirical, or product statement.
\end{boundary}
